\appendix
\section{Toolkit: the five facts we use from outside graph theory}\label{app:toolkit}

Everything in this appendix is standard, and none of it is new. It is collected here, with proofs,
because the rewrite is addressed to a reader who has had one course in graph theory and none in
linear algebra beyond the rank of a matrix. The main text uses \cref{thm:konig,lem:minor-matching}
repeatedly; \cref{thm:farkas} is used exactly once, in \cref{sec:weighted}, and a reader willing to
take the vertex-weighted corollary on trust may skip it.

\subsection{Matchings, covers, and K\H{o}nig's theorem}\label{app:konig}

Let $\Gamma$ be a bipartite graph with parts $A$ and $B$. A \emph{matching} is a set of edges no two
of which share a vertex; $\nu(\Gamma)$ is the largest size of a matching. A \emph{vertex cover} is a
set $K\subseteq A\cup B$ meeting every edge.

\begin{theorem}[K\H{o}nig]\label{thm:konig}
In a bipartite graph $\Gamma$, the maximum size of a matching equals the minimum size of a vertex
cover. Moreover, if $\mathcal M$ is a maximum matching and $A_0$ is the set of vertices of $A$
unmatched by $\mathcal M$, then the set $K$ constructed below is a vertex cover with
$\abs K=\abs{\mathcal M}$ that contains exactly one endpoint of every edge of $\mathcal M$ and no
unmatched vertex.
\end{theorem}

\begin{proof}
A cover has at least one vertex from each edge of a matching, so $\abs K\ge\nu(\Gamma)$ for every
cover $K$; it suffices to exhibit a cover of size $\nu(\Gamma)$. Let $\mathcal M$ be a maximum
matching. Starting from all unmatched vertices of $A$, follow edges \emph{outside} $\mathcal M$ from
$A$ to $B$ and edges \emph{in} $\mathcal M$ from $B$ to $A$; let $W_A\subseteq A$ and $W_B\subseteq B$
be the sets of vertices reached.

\emph{No unmatched vertex of $B$ is reached}: an alternating path from an unmatched vertex of $A$ to
an unmatched vertex of $B$ would be augmenting, and augmenting it would enlarge $\mathcal M$.

\emph{The endpoints of a matched edge are either both reached or both unreached}. If the $B$-endpoint
is reached, the matching edge leads to its $A$-endpoint, so that is reached too; conversely a reached
$A$-vertex that is matched was reached along its matching edge from its partner.

Put $K=(A\setminus W_A)\cup W_B$. It covers every edge: an edge $ab$ with $a\in W_A$ and
$b\notin W_B$ would be neither outside $\mathcal M$ (then $b$ would be reached) nor in $\mathcal M$
(since matched endpoints share their status). It contains exactly one endpoint of each matched edge
and no unmatched vertex, because every unmatched vertex of $A$ lies in $W_A$ and, by the first
paragraph, no unmatched vertex of $B$ lies in $W_B$. Hence
$\abs K=\abs{\mathcal M}=\nu(\Gamma)$.
\end{proof}

\subsection{Ranks}\label{app:ranks}

For a real matrix $M$ with $c$ columns, \emph{rank--nullity} is the statement
$\dim\ker M=c-\rank M$: the rank is the dimension of the image and the domain splits as the kernel
direct sum a complement.

\begin{lemma}[Sylvester's rank inequality]\label{lem:sylvester}
For real matrices $X$ (with $c$ columns) and $Y$ (with $c$ rows),
\[
 \rank(XY)\ \ge\ \rank X-\dim\ker Y .
\]
\end{lemma}

\begin{proof}
The map $Y$ sends $\im X$ onto $\im(XY)$, and its kernel on $\im X$ is $\im X\cap\ker Y$. A linear
map on a space of dimension $d$ with kernel of dimension $e$ has image of dimension $d-e$, so
$\rank(XY)=\dim\im X-\dim(\im X\cap\ker Y)\ge\rank X-\dim\ker Y$.
\end{proof}

\begin{lemma}[A nonzero minor exhibits a matching]\label{lem:minor-matching}
Let $T\subseteq I\times J$ be finite and let $M$ be a real matrix supported on $T$, that is, with
$M_{ij}=0$ whenever $(i,j)\notin T$. Then
\[
 \rank M\ \le\ \nu(T),
\]
where $\nu(T)$ is the largest size of a set of entries of $T$ no two in the same row or column.
In particular, if $\mathcal M$ is a maximum matching in the support $T$, then every matrix supported
on $T$ has rank at most $\abs{\mathcal M}$, and equality is attained by some matrix supported
on $T$.
\end{lemma}

\begin{proof}
Let $r=\rank M$, so $M$ has an $r\times r$ submatrix $\widehat M$ with $\det\widehat M\ne0$. Expanding
the determinant,
\[
 \det\widehat M=\sum_{\sigma}\operatorname{sgn}(\sigma)\prod_{i}\widehat M_{i\sigma(i)},
\]
and every product in the sum is a product of entries of $T$ with distinct rows and distinct columns,
i.e.\ of a matching of size $r$ in $T$; such a matching exists as soon as one product is nonzero, and
one is nonzero because the sum is. Hence $r\le\nu(T)$.

For the last sentence: the identity matrix on a matching of size $\nu(T)$ (entries $1$ on the matched
pairs, $0$ elsewhere) is supported on $T$ and has rank $\nu(T)$.
\end{proof}

\begin{remark}[What we do and do not need]\label{rem:no-edmonds}
The classical statement is that the \emph{maximum rank} over matrices supported on $T$ equals
$\nu(T)$; the second half of \cref{lem:minor-matching} is that statement, in the easy direction,
and we use only this. In the proof of the matching bound (\cref{sec:matching}) the two inequalities
are used as follows: an upper bound $\rank L\le\ell$ for \emph{every} assignment of the variables,
because the support of $L$ contains no matching larger than $\ell$ (\cref{lem:minor-matching},
first part); and the existence of \emph{one} assignment with $\rank L=\ell$, exhibited by an explicit
minor whose determinant is a nonzero polynomial (\cref{sec:matching}, Step~5). No appeal to the
hard direction is needed, and in particular this document does not rely on the
Edmonds/Egerv\'ary rank formula.
\end{remark}

\subsection{The Farkas alternative}\label{app:farkas}

\begin{theorem}[Farkas, in the form we use]\label{thm:farkas}
For a real matrix $M$ with $n$ columns, exactly one of the following holds:
\begin{enumerate}[label=\textup{(\roman*)},leftmargin=1.6em,itemsep=1pt,topsep=1pt]
\item there is $\eta\in\R^n$ with $\eta\ge0$, $\eta\ne0$ and $M\eta\ge0$;
\item there is $\xi\ge0$ with $M^{\mathsf T}\xi<0$ in every coordinate.
\end{enumerate}
\end{theorem}

\begin{proof}[Sketch]
Apply Farkas' lemma in the form: exactly one of
$\{\xi\ge0:\ A\xi\ge b\}\ne\varnothing$ and $\{\eta\ge0:\ A^{\mathsf T}\eta\le0,\ b^{\mathsf T}\eta>0\}\ne\varnothing$
holds. Take $A=-M^{\mathsf T}$ and $b=\mathbf 1$. The first alternative becomes
$\xi\ge0$, $M^{\mathsf T}\xi\le-\mathbf 1$, which is equivalent to (ii) by rescaling; the second
becomes $\eta\ge0$, $M\eta\ge0$, $\mathbf 1^{\mathsf T}\eta>0$, which is (i).
\end{proof}

\begin{remark}[Where each external input is used]\label{rem:inputs}
For the main theorem, the only inputs from outside the definitions are \cref{thm:konig} and
\cref{lem:sylvester,lem:minor-matching} (all proved above, all elementary). The weighted forms of
\cref{sec:weighted} additionally use \cref{thm:farkas} for part (ii) and Seacrest's
equivalence~\cite{Seacrest2015} for part (iii); \cref{sec:cycles} uses nothing beyond
\cref{thm:main}.
\end{remark}
